% sections/biosketch.tex
% Biographical Sketches for PI and Co-PI (NIH Biosketch format)

\section*{Biographical Sketches}

% ============================================================================
\subsection*{Jeffrey R.\ Hill --- Principal Investigator}
% ============================================================================

\subsubsection*{A.\ Personal Statement}

I am an Associate Professor of Mechanical Engineering at Brigham Young
University with expertise in tensegrity structures, smart materials, and
experimental mechanics.  My research group has recently published on tensegrity
robot locomotion, cable tension interactions, and 3D-printed tension networks,
providing direct expertise relevant to this proposal.  I have a strong record
of mentoring undergraduate and graduate students through hands-on research,
which aligns with the goals of the Mentored Research Grant program.

\begin{enumerate}[nosep]
    \item Layer B, Denning H, Hill J. Control and locomotion of tensegrity robots through manipulation of the center of mass. \emph{Robotica}. 2024. \url{https://doi.org/10.1017/S0263574724001139}
    \item Masmeijer T, Swain C, Hill J, Habtour E. Programming tension in 3D printed networks inspired by spiderwebs. \emph{Materials \& Design}. 2025. \url{https://doi.org/10.1016/j.matdes.2025.115281}
\end{enumerate}

\subsubsection*{B.\ Positions, Scientific Appointments, and Honors}

\textbf{Education/Training}

\begin{tabular}{@{}llll@{}}
\textbf{Institution} & \textbf{Degree} & \textbf{Completion Date} & \textbf{Field of Study} \\
\hline
Pennsylvania State University & Ph.D. & 2011 & Mechanical Engineering \\
Brigham Young University & M.S. & 2005 & Mechanical Engineering \\
Brigham Young University & B.S. & 2004 & Mechanical Engineering \\
\end{tabular}

\textbf{Positions and Employment}

\begin{itemize}[nosep]
    \item 2021--present: Associate Professor, Dept.\ of Mechanical Engineering, Brigham Young University
    \item 2016--2021: Principal Member of the Technical Staff, Sandia National Laboratories
    \item 2011--2016: Senior Member of the Technical Staff, Sandia National Laboratories
    \item 2008--2011: Visiting Research Investigator, University of Michigan
\end{itemize}

\subsubsection*{C.\ Contributions to Science}

\textbf{1.\ Tensegrity structures and robots.}
My group investigates design, construction, and control of tensegrity systems.
We have demonstrated simplified construction methods, characterized cable
tension interactions, and developed locomotion strategies for tensegrity
robots.

\begin{enumerate}[nosep]
    \item Layer B, Denning H, Hill J. Control and locomotion of tensegrity robots through manipulation of the center of mass. \emph{Robotica}. 2024. \url{https://doi.org/10.1017/S0263574724001139}
    \item Denning H, Bullinger A, Hill J. Tuning Tensegrity: Simplified Construction and Cable Tension Interactions. \emph{ASME SMASIS 2024}. \url{https://doi.org/10.1115/SMASIS2024-140381}
    \item Brown A, Swain C, Usevitch N, Hill J. Development of a Continuous Cable Tensegrity. \emph{ASME SMASIS 2024}. \url{https://doi.org/10.1115/SMASIS2024-140185}
\end{enumerate}

\textbf{2.\ Advanced manufacturing and structural design.}
I have contributed to 3D-printed tension networks inspired by spiderwebs and
encapsulation techniques for protecting electronic components, demonstrating
expertise in additive manufacturing and structural characterization.

\begin{enumerate}[nosep]
    \item Masmeijer T, Swain C, Hill J, Habtour E. Programming tension in 3D printed networks inspired by spiderwebs. \emph{Materials \& Design}. 2025. \url{https://doi.org/10.1016/j.matdes.2025.115281}
    \item Hill J, Chen A, Wilson N, Boll C, O'Bannon M, Trentman D. Microbead Encapsulation for Protection of Electronic Components. \emph{J.\ Electronic Packaging}. 2025;147(021007). \url{https://doi.org/10.1115/1.4067650}
\end{enumerate}

\textbf{3.\ Smart materials and structural optimization.}
My earlier work focused on actuator grouping optimization and piezoelectric
fiber composite analysis for adaptive structures, establishing a foundation in
computational and experimental approaches to structural systems.

\begin{enumerate}[nosep]
    \item Hill J, Wang K. Development of the En Masse Elimination Algorithm for Actuator Grouping Optimization. \emph{J.\ Intelligent Material Systems and Structures}. 2011;22(11):1203--1212. \url{https://doi.org/10.1177/1045389X11417651}
    \item Kim J, Hill J, Wang K. An asymptotic approach for the analysis of piezoelectric fiber composite beams. \emph{Smart Materials and Structures}. 2011;20(2):025023.
    \item Hill J, Wang K, Fang H, Quijano U. Actuator grouping optimization on flexible space reflectors. \emph{SPIE Active and Passive Smart Structures}. 2011. \url{https://doi.org/10.1117/12.880086}
\end{enumerate}

\vspace{2em}

% ============================================================================
\subsection*{Sterling G.\ Baird --- Co-Principal Investigator}
% ============================================================================

\subsubsection*{A.\ Personal Statement}

I am an Assistant Professor of Mechanical Engineering at Brigham Young
University specializing in Bayesian optimization, self-driving laboratories,
and data-driven materials discovery.  My work on the Honegumi interface and
high-dimensional Bayesian optimization directly supports the optimization
component of this proposal.  I have extensive experience mentoring
undergraduate researchers and developing accessible open-source tools that
lower barriers to adopting advanced optimization methods.

\begin{enumerate}[nosep]
    \item Baird SG, Falkowski AR, Sparks TD. Honegumi: An Interface for Accelerating the Adoption of Bayesian Optimization in the Experimental Sciences. 2025. \url{https://doi.org/10.48550/arXiv.2502.06815}
    \item Tom G, Schmid SP, Baird SG, et al. Self-Driving Laboratories for Chemistry and Materials Science. \emph{Chemical Reviews}. 2024;124(16):9633--9732. \url{https://doi.org/10.1021/acs.chemrev.4c00055}
\end{enumerate}

\subsubsection*{B.\ Positions, Scientific Appointments, and Honors}

\textbf{Education/Training}

\begin{tabular}{@{}llll@{}}
\textbf{Institution} & \textbf{Degree} & \textbf{Completion Date} & \textbf{Field of Study} \\
\hline
University of Utah & Ph.D. & 2023 & Materials Science and Engineering \\
Brigham Young University & M.S. & 2021 & Mechanical Engineering \\
Brigham Young University & B.S. & 2018 & Applied Physics \\
\end{tabular}

\textbf{Positions and Employment}

\begin{itemize}[nosep]
    \item 2025--present: Assistant Professor, Dept.\ of Mechanical Engineering, Brigham Young University
    \item 2025: Principal Investigator and Staff Scientist, Acceleration Consortium, University of Toronto
    \item 2023--2025: Director of Training and Programs, Acceleration Consortium, University of Toronto
\end{itemize}

\textbf{Honors}
\begin{itemize}[nosep]
    \item Guest Editor, \emph{npj Computational Materials}
    \item Journal reviewer for \emph{Nature Communications}, \emph{Digital Discovery}, \emph{JOSS}
\end{itemize}

\subsubsection*{C.\ Contributions to Science}

\textbf{1.\ Bayesian optimization for materials and experimental sciences.}
I developed Honegumi, an interactive interface for generating Bayesian
optimization scripts, and demonstrated high-dimensional BO applied to
materials property prediction and formulation optimization.

\begin{enumerate}[nosep]
    \item Baird SG, Falkowski AR, Sparks TD. Honegumi: An Interface for Accelerating the Adoption of Bayesian Optimization in the Experimental Sciences. 2025. \url{https://doi.org/10.48550/arXiv.2502.06815}
    \item Baird SG, Liu M, Sparks TD. High-Dimensional Bayesian Optimization of 23 Hyperparameters over 100 Iterations for an Attention-Based Network to Predict Materials Property. \emph{Computational Materials Science}. 2022;211:111505. \url{https://doi.org/10.1016/j.commatsci.2022.111505}
    \item Baird SG, Hall JR, Sparks TD. Compactness Matters: Improving Bayesian Optimization Efficiency of Materials Formulations through Invariant Search Spaces. \emph{Computational Materials Science}. 2023;224:112134. \url{https://doi.org/10.1016/j.commatsci.2023.112134}
    \item Baird SG, et al. Bayesian Optimization Hackathon for Chemistry and Materials. 2025. \url{https://doi.org/10.26434/chemrxiv-2025-dzh5z}
\end{enumerate}

\textbf{2.\ Self-driving laboratories.}
I co-authored the definitive review of self-driving labs for chemistry and
materials science and developed a low-cost ``Hello World'' demonstration
platform, making autonomous experimentation accessible to a broad audience.

\begin{enumerate}[nosep]
    \item Tom G, Schmid SP, Baird SG, et al. Self-Driving Laboratories for Chemistry and Materials Science. \emph{Chemical Reviews}. 2024;124(16):9633--9732. \url{https://doi.org/10.1021/acs.chemrev.4c00055}
    \item Baird SG, Sparks TD. Building a ``Hello World'' for Self-Driving Labs: The Closed-loop Spectroscopy Lab Light-mixing Demo. \emph{STAR Protocols}. 2023;4(2):102329. \url{https://doi.org/10.1016/j.xpro.2023.102329}
    \item Lo S, Baird SG, et al. Review of Low-Cost Self-Driving Laboratories in Chemistry and Materials Science: The ``Frugal Twin'' Concept. \emph{Digital Discovery}. 2024;3(5):842--868. \url{https://doi.org/10.1039/D3DD00223C}
\end{enumerate}

\textbf{3.\ Community building and open-source tools for data-driven research.}
I organized a Bayesian optimization hackathon (120 participants, 62
organizations, 14 countries) and developed open-source tools that lower the
barrier to adopting data-driven methods in experimental research.  I have
mentored 8 undergraduate research assistants at BYU and 12+ staff and students
at the Acceleration Consortium.

\begin{enumerate}[nosep]
    \item Baird SG, et al. Bayesian Optimization Hackathon for Chemistry and Materials. 2025. \url{https://doi.org/10.26434/chemrxiv-2025-dzh5z}
    \item Seegmiller CC, Baird SG, Sayeed HM, Sparks TD. Discovering Chemically Novel, High-Temperature Superconductors. \emph{Computational Materials Science}. 2023.
\end{enumerate}
